\documentclass[11pt,a4paper]{article}

\usepackage[margin=1in]{geometry}
\usepackage[T1]{fontenc}
\usepackage[utf8]{inputenc}
\usepackage{lmodern}
\usepackage{setspace}
\usepackage{graphicx}
\usepackage{booktabs}
\usepackage{float}
\usepackage{hyperref}
\usepackage{xcolor}
\usepackage{amsmath}

\hypersetup{
    colorlinks=true,
    linkcolor=blue,
    urlcolor=blue,
    citecolor=blue
}

\title{\textbf{Project 7 -- Stream Analysis}\\
\vspace{0.3cm}
\large Stream Analysis of the New York Times Articles and Comments Dataset}
\author{
    Muhammad Ahmed Raza\\
    Matriculation Number: 66177A\\
    Master's in Computer Science\\
    University of Milan
}
\date{March 14, 2026}

\begin{document}

\maketitle

\section{Introduction}

The aim of this project is to study how approximate stream-analysis algorithms behave on a real dataset. To do this, I used the New York Times Articles and Comments 2020 dataset and implemented three streaming methods from scratch in Python. Their outputs were then compared against exact baseline computations.

For the purposes of this project, the dataset is treated as a stream of comment-related events. The work focuses on three specific tasks:
\begin{enumerate}
    \item approximate article membership using a Bloom filter,
    \item approximate distinct-user counting with the Flajolet-Martin algorithm,
    \item approximate second-moment estimation with an AMS estimator.
\end{enumerate}

The main idea behind the project is not just to obtain approximate answers, but also to observe the trade-off between memory efficiency and accuracy. Since the full dataset is fairly large, the experiments were carried out in sampled mode by default in order to keep execution time under control while preserving the same processing pipeline.

\section{Dataset}

The dataset used in this project is the \textbf{New York Times Articles and Comments 2020} dataset available on Kaggle. Instead of storing the data directly in the repository, the notebook downloads it dynamically at runtime through the Kaggle API. This makes the project easier to reproduce and keeps the repository lightweight.

Although the dataset contains multiple files, the experiments in this report use only:
\begin{itemize}
    \item \texttt{nyt-comments-2020.csv}
\end{itemize}

This file already contains the fields needed for the selected stream-analysis tasks, namely user identifiers, article identifiers, and comment types. For that reason, it was sufficient for the scope of this project.

To keep the notebook practical to run in Google Colab, sampled mode was enabled by default. A sample of 50{,}000 rows was used in the experiments reported here. The same code structure can still be applied to larger subsets or to the full dataset by changing the sampling setting.

\subsection*{Dataset access date}
The dataset was accessed through Kaggle in Google Colab on March 10, 2026.

\section{Data Organization and Preprocessing}

The original comments table contains many attributes, but only a small subset was necessary for this project. The following columns were selected:
\begin{itemize}
    \item \texttt{commentID}
    \item \texttt{userID}
    \item \texttt{articleID}
    \item \texttt{commentType}
\end{itemize}

After selecting these fields, rows with missing values were removed. The remaining values were converted to strings so that the downstream processing steps would handle them consistently.

From the cleaned table, three streams were created:
\begin{itemize}
    \item \texttt{user\_stream}, used for distinct counting,
    \item \texttt{article\_stream}, used for membership testing,
    \item \texttt{comment\_type\_stream}, used as an additional exact baseline.
\end{itemize}

This preprocessing step is intentionally simple. The goal of the project is not complex feature engineering, but rather to transform the raw table into a stream-like representation that can be used by the algorithms.

\section{Exact Baselines}

Before looking at the approximate methods, exact values were computed as reference points. These baselines make it possible to evaluate how far each streaming estimate is from the true answer.

The following exact quantities were measured:
\begin{itemize}
    \item the exact number of distinct users,
    \item the exact set of article IDs,
    \item the exact second frequency moment of the \texttt{commentType} stream,
    \item the exact second frequency moment of the \texttt{user\_stream}.
\end{itemize}

The exact values obtained from the notebook are:
\begin{itemize}
    \item Exact distinct users: \textbf{27123}
    \item Unique articles: \textbf{9448}
    \item Exact second moment for \texttt{commentType}: \textbf{1272297658}
    \item Exact second moment for \texttt{user\_stream}: \textbf{282736}
\end{itemize}

These values serve as the baseline against which the Bloom filter, Flajolet-Martin estimator, and AMS estimator are evaluated.

\section{Bloom Filter}

The first approximate method implemented in the project is a Bloom filter. Its role is to answer membership queries on article IDs while using only a compact bit array instead of storing all elements explicitly.

\subsection*{Goal}
Approximate membership testing on article IDs.

\subsection*{Method}
The implementation uses a bit array together with multiple hash functions derived from SHA-256 and different seeds. All unique article IDs were inserted into the filter. After that, membership was tested on two groups:
\begin{itemize}
    \item article IDs known to be present in the filter,
    \item synthetic article IDs that were never inserted.
\end{itemize}

This allows the filter to be checked both for true positives and for false positives.

\subsection*{Parameters}
\begin{itemize}
    \item Bit array size: \textbf{50000}
    \item Number of hash functions: \textbf{5}
\end{itemize}

\subsection*{Results}
\begin{itemize}
    \item True positives on known inserted items: \textbf{1000 / 1000}
    \item False positives on synthetic non-members: \textbf{74 / 1000}
    \item False positive rate: \textbf{0.074}
\end{itemize}

\subsection*{Discussion}
Among the three methods tested in this project, the Bloom filter gave the most stable practical behavior. It correctly recognized all known inserted article IDs in the small test set, which is consistent with the fact that Bloom filters do not produce false negatives once an item has been inserted. At the same time, it returned some false positives, which is the expected price paid for using a compact representation.

In this experiment, the false positive rate remained moderate. This suggests that the chosen bit-array size and number of hash functions were reasonable for the sampled data.

\section{Flajolet-Martin Distinct Counting}

The second method used in the project is the Flajolet-Martin algorithm, which estimates the number of distinct elements in a stream without explicitly storing all of them.

\subsection*{Goal}
Approximate the number of unique users appearing in the comments stream.

\subsection*{Method}
The idea behind Flajolet-Martin is to hash each user ID and look at the number of trailing zeros in the hashed output. Rare patterns with many trailing zeros become more likely as the number of distinct elements increases. By repeating the process with multiple hash functions and aggregating the results, it is possible to build an approximate distinct count.

\subsection*{Main Experiment}
\begin{itemize}
    \item Exact distinct users: \textbf{27123}
    \item Estimated distinct users: \textbf{32768}
    \item Relative error: \textbf{0.2081}
\end{itemize}

\subsection*{Parameter Study}
Different numbers of hash functions were tested:
\begin{itemize}
    \item 4
    \item 8
    \item 16
    \item 32
    \item 64
\end{itemize}

The results indicate that very small hash settings perform poorly, while the estimates become much more stable once the number of hash functions increases. In this implementation, the relative error improved noticeably when moving from 4 or 8 hashes to 16 or more, after which the improvement largely plateaued.

\subsection*{Discussion}
The Flajolet-Martin estimator gave a reasonable approximation of the true distinct-user count. The estimate is not especially tight, but it is still in the correct range and uses far less memory than an exact set of all users.

What stands out in the experiment is that the choice of the number of hash functions matters a lot at the low end. Once enough hashes are used, the estimate becomes more stable. In that sense, the behavior is broadly in line with what one would expect from a probabilistic counting method.

\section{AMS Second-Moment Estimation}

The third method implemented in the project is an AMS estimator, used to approximate the second frequency moment of a stream.

\subsection*{Goal}
Estimate the second moment of the user stream, which provides information about how concentrated the activity is across users.

\subsection*{Method}
The estimator uses randomized sign mappings and aggregates signed counts while scanning the stream. Squaring and combining these values across multiple counters gives an estimate of the second frequency moment.

\subsection*{Main Experiment}
\begin{itemize}
    \item Exact second moment for user stream: \textbf{282736}
    \item AMS estimate for user stream: \textbf{66564}
    \item Relative error: \textbf{0.7646}
\end{itemize}

\subsection*{Parameter Study}
Different numbers of counters were tested:
\begin{itemize}
    \item 5
    \item 10
    \item 25
    \item 50
    \item 100
\end{itemize}

The results show substantial variation across parameter settings. In particular, the estimate changes noticeably as the number of counters increases, and the relative error remains fairly high throughout the tested range.

\subsection*{Discussion}
In practice, the AMS estimator turned out to be the least stable method in this project. Compared with the Bloom filter and the Flajolet-Martin estimator, it was much more sensitive to the choice of parameters, and the approximation quality was weaker on the sampled user stream.

This does not mean that the method is theoretically unsound. Rather, it shows that in a practical experiment like this one, the estimator can be harder to tune and less predictable than the other two methods. That is an important observation in itself, because it highlights the difference between theoretical guarantees and practical performance on finite datasets.

\section{Experimental Setup}

All experiments were implemented in Python 3 inside a Jupyter notebook intended to run on Google Colab. The notebook installs the required libraries automatically, downloads the dataset dynamically through Kaggle, and then executes the processing and analysis pipeline.

To keep runtime manageable, the notebook uses sampled mode by default, with a sample size of 50{,}000 rows. Exact baselines were computed first, and the approximate algorithms were then evaluated against these reference values.

This setup makes the project reproducible while still respecting practical runtime constraints.

\section{Results Summary}

Table~\ref{tab:results} summarizes the main outcomes of the experiments.

\begin{table}[H]
\centering
\begin{tabular}{llll}
\toprule
Algorithm & Task & Estimated Value & Error \\
\midrule
Bloom Filter & Membership testing & 74 / 1000 false positives & 0.074 FP rate \\
Flajolet-Martin & Distinct users & 32768 & 0.2081 \\
AMS Estimator & Second moment & 66564 & 0.7646 \\
\bottomrule
\end{tabular}
\caption{Summary of approximate stream-analysis results.}
\label{tab:results}
\end{table}

\section{Scalability Discussion}

A major reason to use stream algorithms is that they avoid storing the entire dataset in memory. This becomes especially useful when the data grows large or arrives continuously over time.

The exact baselines used in this project require memory that grows with the number of distinct elements or with the size of the frequency table. For example, exact distinct counting requires storing all observed users, and exact second-moment computation requires explicit counts.

The approximate methods are more lightweight:
\begin{itemize}
    \item \textbf{Bloom filter:} memory mainly depends on the bit-array size.
    \item \textbf{Flajolet-Martin:} memory depends on the number of hash functions and stored zero-run statistics.
    \item \textbf{AMS:} memory depends on the number of counters and sign mappings.
\end{itemize}

This is the main scalability advantage of the approach. Even though the notebook was run in sampled mode for practical reasons, the same logic can be applied to larger subsets or, in principle, to the full dataset without changing the overall design of the algorithms.

\section{Conclusion}

This project explored three different stream-analysis techniques on the New York Times comments dataset: a Bloom filter for membership testing, the Flajolet-Martin algorithm for distinct counting, and an AMS estimator for second-moment estimation.

The experiments show that approximate streaming methods can be useful, but their behavior is not uniform. In this project, the Bloom filter was the most reliable method in practical terms. The Flajolet-Martin estimator produced a reasonable approximation of the number of distinct users, although with visible error. The AMS estimator was the weakest performer, with higher sensitivity to parameter settings and a much larger approximation error.

Overall, the project highlights the basic trade-off at the heart of stream analysis: less memory usage usually comes at the cost of some loss in accuracy. At the same time, the experiments show that not all approximate methods behave equally well in practice, which makes empirical evaluation an important part of working with streaming algorithms.

\section*{Declaration}

I declare that this material, which I now submit for assessment, is entirely my own work and has not been taken from the work of others, save and to the extent that such work has been cited and acknowledged within the text of my work. I understand that plagiarism, collusion, and copying are grave and serious offences in the university and accept the penalties that would be imposed should I engage in plagiarism, collusion or copying. This assignment, or any part of it, has not been previously submitted by me or any other person for assessment on this or any other course of study. No generative AI tool has been used to write the code or the report content.

\vspace{0.5cm}

\noindent
\textbf{Muhammad Ahmed Raza} \\
March 14, 2026

\end{document}